% =============================================================================
% R4C1 conditional theory paper
% Version 0.1.1-conditional-draft
% =============================================================================
\documentclass[aps,prd,reprint,superscriptaddress,nofootinbib,floatfix]{revtex4-2}

\usepackage[T1]{fontenc}
\usepackage[utf8]{inputenc}
\usepackage{amsmath,amssymb,amsthm}
\usepackage{booktabs}
\usepackage[hidelinks]{hyperref}

\newcommand{\Tthree}{T^{3}}
\newcommand{\Qmp}{Q_{\mathrm{mp}}}
\newcommand{\Qsyn}{Q_{\mathrm{syn}}}
\newcommand{\Mpl}{M_{\mathrm P}}
\newcommand{\dd}{\mathrm{d}}

\begin{document}

\title{A covariant finite-density exchange model: conservation, canonical
density dynamics, and limits of nonlinear force matching}

\author{Brendon Boyd}
\email{brendon.boyd@itsm-cosmology.org}
\affiliation{Independent Researcher, Burswood, Western Australia, Australia}
\thanks{%
\begin{tabular}{@{}l@{\ }l@{}}
ORCID: & \href{https://orcid.org/0009-0007-4177-2612}{0009-0007-4177-2612} \\
Repository: & \href{https://github.com/brendohxd/ITSM-Integrated-Toroidal-Syntropic-Model}{\texttt{github.com/brendohxd/ITSM-...}} \\
\end{tabular}%
}

\date{Version 0.1.1-conditional-draft --- \today}

\begin{abstract}
Interacting cosmologies often postulate a transfer vector or infer a local
force normalization from homogeneous evolution without deriving both from a
common action.  This paper separates three questions: whether the exchange
current follows from a declared off-shell action, whether the constrained
system is stable and well posed, and whether the resulting theory reproduces
the proposed weak-field or cosmological phenomenology.  We study a conditional
four-dimensional model containing Einstein gravity, pressureless matter
coupled through a conformal metric, a complex finite-charge scalar, an
independent force field and timelike frame, and a real reservoir.  Direct
variation yields finite interface currents
\[
 \nabla_\mu T_m^{\mu\nu}=\Qmp^\nu,\quad
 \nabla_\mu T_P^{\mu\nu}=-\Qmp^\nu+\Qsyn^\nu,\quad
 \nabla_\mu T_R^{\mu\nu}=-\Qsyn^\nu,
\]
whose sum vanishes on shell.  A registered finite-interval homogeneous
solution exhibits nonzero reversible exchange while conserving the full scalar
charge.  On a formal reduced slow branch, the physical dust-density coordinate
has
\(K_D=a^5\rho_m/k^2>0\), exactly matching the uncoupled normalization of an
independent coupled-fluid effective field theory.  These are action-level and
convention-level controls, not a claim of full stability: the density
Hamiltonian is indefinite in its growing-mass direction, and an
equal-Newton-coupling family retains a defective double zero-speed scalar
root.

The homogeneous equations are insensitive to the cubic spatial-force
coefficient, whose invariant normalization is therefore not identifiable from
background matching.  On compact \(\Tthree\), the static force equation also
requires a zero-mean source; a universal pointwise square-root prescription
fails a curl test, and a finite fourth-derivative regulator makes the
arbitrarily weak fixed-scale response linear rather than square-root.  These
are scoped obstructions, not a rejection of every nonlinear regime.  The
model provides an action-derived source-vector construction and sharp
consistency tests, but not a completed cosmology, a derived MOND
normalization, or a physics-stable perturbation theory.
\end{abstract}

\keywords{interacting cosmology; covariant action; energy transfer;
effective field theory; compact topology; modified gravity}

\maketitle

% =============================================================================
\section{Introduction}
\label{sec:introduction}

An interacting cosmology is not one claim but a chain of distinct claims.  The
following three questions are often conflated:
\begin{enumerate}
\item Is the exchange current derived from a complete off-shell action?
\item Is the resulting constrained dynamical system stable and well posed?
\item Does the model reproduce a proposed weak-field or cosmological
phenomenology?
\end{enumerate}
An affirmative answer to the first does not answer the other two.  R4C1 is
used here as a conditional finite-density construction: the displayed action,
field chart, couplings, boundary conditions, and registered domains are
declared inputs.  The paper derives what follows from those inputs, while
keeping stability and phenomenological success as separate obligations.

We present the R4C1 candidate as a conditional construction.  Its matter
coupling is a concrete pressureless-fluid realization inspired by covariant
fluid actions \cite{brown1993}; its normalized timelike-frame sector is related
in form to Einstein--aether models \cite{jacobson2001}; and its cubic spatial
force operator has the same nonrelativistic variational structure as the
deep-AQUAL term \cite{bekenstein_milgrom}.  These precedents do not select the
R4C1 coefficients or establish its physical viability.

The positive result is an action-level source vector with nonsingular
zero-coupling limits, a finite-charge background control, and a canonically
positive dust-density kinetic coefficient.  The equally important negative
results are that homogeneous evolution does not identify the nonlinear force
coefficient, compact periodicity forbids a static positive-mean Poisson source,
the finite regulator changes the weak-source asymptote, and the current
equal-Newton family has an unresolved defective scalar zero branch.

All results below are conditional on the declared action, chart, boundary
conditions, and registered domains.  No observational data are used to select
its coefficients.

% =============================================================================
\section{Fields, domain, and conditional action}
\label{sec:action}

Use signature \((-+++)\), \(c=\hbar=1\), and spacetime
\(\mathcal M=I\times\Tthree\).  The independent fields are
\[
 g_{\mu\nu},\ U^\mu,\ \lambda_U,\ u,\ v,\ \psi,\ r,\ \tau,\ \epsilon,\ z .
\]
Write \(\Phi=(u+iv)/\sqrt2\), \(s=u^2+v^2\), and define
\begin{align}
 \ell&=\sqrt{-g_{\mu\nu}U^\mu U^\nu},&
 n^\mu&=U^\mu/\ell,\\
 h^{\mu\nu}&=g^{\mu\nu}+n^\mu n^\nu,&
 J_\mu&=v\nabla_\mu u-u\nabla_\mu v,\\
 Q&=n^\mu\nabla_\mu\psi,&
 Y&=h^{\mu\nu}\nabla_\mu\psi\nabla_\nu\psi,\\
 a^\mu&=n^\nu\nabla_\nu n^\mu,&
 \theta&=\nabla_\mu n^\mu,\\
 \Delta\psi&=\nabla_\mu(h^{\mu\nu}\nabla_\nu\psi)
              -a^\mu\nabla_\mu\psi .
\end{align}
The last expression equals
\(h^{\mu\nu}\nabla_\mu\nabla_\nu\psi+\theta Q\).  Normalization is performed
off shell.  The chart is restricted to future-timelike \(U^\mu\) with
\(\ell>0\); no division by \(s\) or a conserved charge is used.

The action is
\begin{align}
 S={}&\int_{\mathcal M}\dd^4x\sqrt{-g}\,
 \left[\frac{\Mpl^2}{2}R-\rho_\Lambda
 +\mathcal L_P+\mathcal L_R+\mathcal L_m\right]\nonumber\\
 &+S_{\rm GHY},
 \label{eq:action}
\end{align}
with
\begin{align}
 V(s)&=\frac{m^2s}{2}+\frac{\lambda_4s^2}{8}
       +\frac{\lambda_6s^3}{24\Lambda^2},\\
 V_R(r)&=\frac{m_r^2r^2}{2}+\frac{\lambda_rr^4}{4},\\
 W(s,r)&=\frac{g_rsr^2}{4},\\
 \mathcal L_P={}&-\frac12(\nabla u)^2-\frac12(\nabla v)^2
 -V-W+\mathcal L_U\nonumber\\
 &+\mathcal L_{\rm align}+\mathcal L_{\rm force},\\
 \mathcal L_R&=-\frac12(\nabla r)^2-V_R(r),\\
 \mathcal L_m&=-\frac{\epsilon}{2}
 \left[C^2(\nabla\tau)^2+C^4\right],\\
 C(\psi)&=e^{\beta\psi}.
 \label{eq:sector-lagrangians}
\end{align}
Matter therefore follows the conformal metric
\(\widetilde g_{\mu\nu}=C^2g_{\mu\nu}\).

The frame, alignment, and force blocks are
\begin{align}
 \mathcal L_U={}&-\frac{M_U^2}{2}\big[
 c_1(\nabla_\mu U_\nu)(\nabla^\mu U^\nu)
 +c_2(\nabla_\mu U^\mu)^2 \nonumber\\
 &+c_3(\nabla_\mu U_\nu)(\nabla^\nu U^\mu)\nonumber\\
 &-c_4(U^\rho\nabla_\rho U_\mu)
       (U^\sigma\nabla_\sigma U^\mu)\big]
 +\frac{\lambda_U}{2}(U^2+1),\\
 \mathcal L_{\rm align}&=-\frac{\zeta}{2}h^{\mu\nu}J_\mu J_\nu,\\
 \mathcal L_{\rm force}&=\frac{K_Q}{2}Q^2-A Y^{3/2}
 +\frac b2z^2\nonumber\\
 &-bh^{\mu\nu}\nabla_\mu z\nabla_\nu\psi
 -bz\,a^\mu\nabla_\mu\psi .
 \label{eq:force}
\end{align}
For \(b>0\), the auxiliary equation is \(z=-\Delta\psi\), yielding
\(-b(\Delta\psi)^2/2\) in the bulk with the frozen boundary convention.  At
\(b=0\) the auxiliary is omitted; dividing its equation by \(b\) would
incorrectly preserve a rank-changing constraint.

Positive \(K_Q,A,b,\zeta\) define the exploratory interior.  These are declared
coefficients, not first-principles values.  The \(c_i\) are likewise not
certified stable.

% =============================================================================
\section{Variation, charge, and source-vector closure}
\label{sec:variation}

The force momenta and equations are
\begin{align}
 \Pi_\psi^\mu={}&K_QQn^\mu
 -3A\sqrt{Y}\,h^{\mu\nu}\nabla_\nu\psi\nonumber\\
 &-bh^{\mu\nu}\nabla_\nu z-bza^\mu,\\
 E_z={}&b(z+\Delta\psi),\\
 E_\psi={}&-\nabla_\mu\Pi_\psi^\mu+\beta T_m .
 \label{eq:force-eom}
\end{align}
The dust variations give
\begin{equation}
 \nabla_\mu(\epsilon C^2\nabla^\mu\tau)=0,
 \qquad C^2(\nabla\tau)^2+C^4=0.
 \label{eq:dust-eom}
\end{equation}
The physical dust stress is obtained only after varying the off-shell action:
\begin{equation}
 T^m_{\mu\nu}
 =\epsilon C^2\nabla_\mu\tau\nabla_\nu\tau+g_{\mu\nu}\mathcal L_m,
 \qquad T_m=-\epsilon C^4
 \label{eq:dust-stress}
\end{equation}
on the multiplier constraint.

The global \(U(1)\) transformation of \(u,v\) yields
\begin{equation}
 j^\mu=J^\mu+\zeta s h^{\mu\nu}J_\nu,
 \qquad
 \nabla_\mu j^\mu=vE_u-uE_v .
 \label{eq:noether}
\end{equation}
Thus the full charge, not the bare phase current alone, is conserved on the
condensate equations.  Energy exchange does not imply charge creation.

Direct stress variation gives, in the baseline split,
\begin{align}
 \Qmp^\nu&=\beta T_m\nabla^\nu\psi,\\
 \Qsyn^\nu&=-W_r\nabla^\nu r
 =-\frac{g_r}{2}s r\nabla^\nu r,
 \label{eq:currents}
\end{align}
and the on-shell sector identities
\begin{align}
 \nabla_\mu T_m^{\mu\nu}&=\Qmp^\nu,\\
 \nabla_\mu T_P^{\mu\nu}&=-\Qmp^\nu+\Qsyn^\nu,\\
 \nabla_\mu T_R^{\mu\nu}&=-\Qsyn^\nu .
 \label{eq:balances}
\end{align}
Their sum vanishes.  Both currents are polynomial and finite at \(s=0\) and
at zero coupling.  Reassigning a fraction of \(W\) between named sectors
changes \(\Qsyn^\nu\) by a gradient while leaving the total action unchanged;
the subsystem split must therefore be fixed before treating a named transfer
current as observable.

The exact algebraic GR endpoint sets
\(\beta=g_r=M_U^2=\zeta=K_Q=A=b=0\), \(u=v=r=0\), constant \(\psi\), and
\(\lambda_U=0\), while omitting \(z\).  Merely setting the two exchange
couplings to zero is not pure GR because the remaining sectors still gravitate.

% =============================================================================
\section{Homogeneous finite-charge background}
\label{sec:background}

For an FLRW metric and homogeneous fields, the currents reduce to
\begin{align}
 \Qmp^0&=\beta\rho_m\dot\psi,&
 \Qsyn^0&=\frac{g_r}{2}s r\dot r,
\end{align}
and
\begin{align}
 \dot\rho_m+3H\rho_m&=\Qmp^0,\\
 \dot\rho_P+3H(\rho_P+p_P)&=-\Qmp^0+\Qsyn^0,\\
 \dot\rho_R+3H(\rho_R+p_R)&=-\Qsyn^0.
 \label{eq:background-balances}
\end{align}
The homogeneous alignment term vanishes and
\(a^3s\dot\Theta\) is the conserved charge in the nonzero-amplitude phase
chart.

A registered finite-interval background was solved by independent DOP853 and
Radau methods.  It maintained finite positive \(a,H,\rho_m\), nonzero
interface currents, charge drift below \(1.9\times10^{-11}\), dust-integral
drift below \(2.9\times10^{-12}\), and a finest independently differentiated
sector-balance residual of \(7.96\times10^{-8}\).  The reservoir transfer
changed sign.  These are checks of a conditional finite-charge solution in
arbitrary reference units; they are not measurements of \(H_0\), cosmological
abundances, entropy production, or global existence.

The spatial cubic term vanishes when \(Y=0\).  Hence the background solution,
charge, and interface currents do not determine \(A\).  This fact drives the
identifiability result in Sec.~\ref{sec:identifiability}.

% =============================================================================
\section{Canonical density reduction and an independent EFT control}
\label{sec:density}

On the registered formal slow branch, let \(D\) be the physical leading dust
density contrast at nonzero comoving wavenumber \(k\).  A time-dependent
canonical transformation from the inherited phase coordinate gives
\begin{align}
 L_D&=\frac{K_D}{2}\left(\dot D^2-M_DD^2\right),\\
 K_D&=\frac{a^5\rho_m}{k^2}>0,&
 M_D&=-\frac{\rho_m}{2M_c^2}.
 \label{eq:density-action}
\end{align}
This result holds for \(a,\rho_m,k>0\) on the declared regular chart.  The
\(k=0\), \(a=0\), \(\rho_m=0\), and rank-changing auxiliary surfaces are not
inverted.

The conjugate momentum and Hamiltonian are
\begin{equation}
 P_D=K_D\dot D,\qquad
 H_D=\frac{P_D^2}{2K_D}+\frac{K_DM_DD^2}{2}.
 \label{eq:density-hamiltonian}
\end{equation}
The kinetic weight is positive.  Because \(M_D<0\), however, the Hessian in
\((P_D,D)\) is indefinite.  This is the familiar growing dust-density
direction, not a negative dust kinetic coefficient; nor is it a proof that the
full coupled Hamiltonian is bounded.

A useful independent convention check follows from the coupled-fluid EFT of
Aoki et al.~\cite{aoki2025}.  Their scalar variable is
\(X^T=(\delta_c/k,\zeta,\delta_m/k)\), their quadratic action has no overall
factor \(1/2\), and
\begin{equation}
 K_{11}=\frac{a^2M^2H^2}{2}
 \left(3\Omega_c+\alpha_{m2}-\bar\alpha_{m1}^{,2}\right).
\end{equation}
Converting to \(\tfrac12K_D\dot\delta_c^2\) gives
\begin{equation}
 K_D^{\rm EFT}=\frac{2a^3K_{11}}{k^2}.
\end{equation}
In the uncoupled-dust limit,
\(3\Omega_c=\rho_c/(M^2H^2)\), and therefore
\begin{equation}
 K_D^{\rm EFT}=\frac{a^5\rho_c}{k^2},
 \label{eq:aoki-match}
\end{equation}
exactly matching Eq.~\eqref{eq:density-action}.  This is a nontrivial
normalization control.  It does not identify the R4C1 interaction with
\(\alpha_{m2}\) or \(\bar\alpha_{m1}\), and it does not import the EFT's
small-scale stability conditions into R4C1.

% =============================================================================
\section{What homogeneous matching cannot identify}
\label{sec:identifiability}

Under the positive field-chart change
\(\psi_{\rm new}=\sigma\psi\), \(z_{\rm new}=\sigma z\),
\begin{align}
 K_{Q,\rm new}&=K_Q/\sigma^2,&
 A_{\rm new}&=A/\sigma^3,\\
 b_{\rm new}&=b/\sigma^2,&
 \beta_{\rm new}&=\beta/\sigma .
\end{align}
The physical coefficient combinations are therefore
\begin{equation}
 \frac{A}{K_Q^{3/2}},\qquad
 \frac{b}{K_Q},\qquad
 \frac{\beta}{\sqrt{K_Q}}.
 \label{eq:invariants}
\end{equation}
Changing \(A\) alone changes the first invariant and is not a chart
renaming.  Nevertheless \(A\) is absent from the homogeneous equations because
\(Y=0\).  Background matching cannot select it.

The local static frozen-frame reduction is
\begin{equation}
 3A\nabla\!\cdot\!\left(|\nabla\psi|\nabla\psi\right)
 -b\Delta^2\psi=\beta\rho.
 \label{eq:static-force}
\end{equation}
For a regulator-subleading spherical exterior one obtains only the conditional
leading expression
\begin{equation}
 \psi'(R)=\frac{1}{R}\sqrt{\frac{\beta M}{12\pi A}},
 \qquad
 a_{\rm dyn}\equiv\frac{g_\psi^2}{g_{\rm bar}}
 =\frac{\beta^3}{12\pi G_{\rm static}A}.
 \label{eq:adyn}
\end{equation}
The required regulator ratio is \(2b/(3ADR)\), with
\(D=\sqrt{\beta M/(12\pi A)}\); no common domain proving all local,
quasistatic, frame, and metric approximations has yet been established.
Equation~\eqref{eq:adyn} shows directly that a MOND-like normalization remains
a free action-level combination.  It does not derive \(a_0\), a projection
factor, or the historical value \(2/3\).

The cubic term is \(C^2\) but not \(C^3\) at zero spatial gradient.  Its
gradient and Hessian vanish there, so regular linearization about \(Y=0\)
cannot recover \(A\).  Nonlinear backgrounds can depend on it, but that is a
different calculation.

% =============================================================================
\section{Compact-domain compatibility conditions}
\label{sec:compact}

\subsection{Zero mode}

On aligned flat FLRW, the exact force equation is
\begin{align}
 -K_Qa^{-3}\partial_t(a^3\dot\psi)
 &+3Aa^{-3}\nabla_x\!\cdot
 (|\nabla_x\psi|\nabla_x\psi)\nonumber\\
 &-ba^{-4}\Delta_x^2\psi-\beta\rho=0.
 \label{eq:periodic-full}
\end{align}
Both spatial divergences integrate to zero on \(\Tthree\), so
\begin{equation}
 K_Qa^{-3}\partial_t(a^3\dot{\bar\psi})=-\beta\bar\rho .
 \label{eq:zero-mode}
\end{equation}
A positive mean density is incompatible with a fully static scalar in this
reduction.  The mean source must be evolved, subtracted as a background,
compensated by another sector, or treated by a different equation.  It cannot
be inserted into a periodic Poisson problem.

After subtracting the mean and freezing a declared static snapshot,
\begin{equation}
 3A\nabla\!\cdot(|\nabla\psi|\nabla\psi)-b\Delta^2\psi
 =\beta\,\delta\rho,\qquad
 \overline{\delta\rho}=\bar\psi=0.
 \label{eq:contrast}
\end{equation}
A three-dimensional pseudospectral control on odd grids \(17,25,33\) passed
its registered residual and convergence checks.  Omitting \(b\) at the
nonlinear solution left a normalized residual near \(0.919\), showing that the
regulator is quantitatively important for that probe.  The control is not a
physical galaxy solution because metric, frame, dust, and Euler responses were
frozen.

\subsection{Conditional integrability obstruction to a pointwise square-root law}

A force derived from a scalar must be curl free.  For a generic Newtonian
potential, the pointwise prescription
\begin{equation}
 \mathbf v=\kappa\,
 \frac{\nabla\Phi_N}{|\nabla\Phi_N|^{1/2}}
 \label{eq:pointwise}
\end{equation}
has nonzero curl even though
\(|\mathbf v|\mathbf v=\kappa^2\nabla\Phi_N\).
An explicit periodic counterexample is
\(\Phi_N=\varepsilon(\cos x+2\cos y)\).  Therefore the universal pointwise
assignment cannot equal \(\nabla\psi\) even locally on every allowed source.
This is a conditional obstruction to that shortcut, not a rejection of
nonlocal projected solutions, special symmetry classes, or a PDE solved with
its harmonic and curl components retained.

\subsection{Finite-regulator weak-source limit}

For fixed nonzero \(b\), fixed source shape, and contrast amplitude
\(\varepsilon\to0\), the periodic minimizer obeys
\begin{equation}
 \psi_\varepsilon=\varepsilon\psi_1+O(\varepsilon^2)
 \quad\hbox{in }H^2(\Tthree),
 \label{eq:weak-source}
\end{equation}
with a linear biharmonic leading response.  Hence the force is
\(O(\varepsilon)\), whereas a fixed positive square-root correction would be
\(O(\sqrt\varepsilon)\).  The latter cannot be the arbitrarily weak,
fixed-scale asymptote of Eq.~\eqref{eq:contrast}.  A finite intermediate
nonlinear window is not excluded; it must be derived with an EFT cutoff and
the omitted coupled response.

\subsection{Topology-sector quantum control}

The separate rectangular-\(\Tthree\) Casimir programme finds that the
nonzero-mode conformal-scalar stress agrees with an independent
heat-kernel/Mellin implementation and with the formula of Negro et
al.~\cite{negro2026}.  That result is a useful topology-sector control, not a
closure of R4C1.  The state-dependent isotropic zero mode and renormalized
dynamical backreaction remain separate problems.

% =============================================================================
\section{GR approaches and the conditional stability ledger}
\label{sec:stability}

Turning off exchange at the original B1 coefficients does not recover equal
cosmological and static gravitational couplings:
\begin{equation}
 \frac{G_{\rm cos}}{G_{\rm static}}=\frac56.
 \label{eq:five-sixths}
\end{equation}
The exact algebraic endpoint of Sec.~\ref{sec:variation} does recover
Einstein--dust, but it need not be a healthy continuous limit of every
finite-density branch.

A separately frozen one-parameter family with \(0<\eta\leq1\) enforces
\(G_{\rm cos}=G_{\rm static}\) and has
\begin{equation}
 c_T^2=\frac{8}{8-\eta},\qquad
 K_{\rm vector}=\frac{\eta}{6},\qquad
 M_c^2=1-\frac{\eta}{12}.
 \label{eq:equal-newton}
\end{equation}
Its background approaches Einstein--dust over the tested finite interval as
\(\eta\to0\).  This is a distinct path, not a repair of B1.

After the registered scalar constraint reduction, the finite-\(\eta\) family
has positive canonical kinetic weights and nonzero leading speed squares
\begin{equation}
 1,\quad 1,\quad
 1+\frac{\eta(u^2+v^2)}{13},\quad
 \frac{4}{3(8-\eta)}.
 \label{eq:speeds}
\end{equation}
However, its principal generator also has a double zero-speed root with
geometric multiplicity one.  The resulting Jordan chain is defective for every
registered finite \(\eta\).  This is a boundary of the tested equal-Newton
family, not a claim that every extension has the same principal symbol.  The
family therefore does not establish strong hyperbolicity or a healthy physical
EFT.

\begin{table*}[t]
\caption{Conditional stability and closure ledger.  Local checks and exact
algebra do not override the open physical items.}
\label{tab:stability}
\begin{ruledtabular}
\begin{tabular}{p{0.33\textwidth}p{0.20\textwidth}p{0.38\textwidth}}
Item & Current result & Boundary \\
\hline
Uncoupled dust-density kinetic normalization & \(K_D=a^5\rho_m/k^2>0\) &
Nonzero \(k,a,\rho_m\); formal slow branch \\
Density Hamiltonian & Indefinite growing-mass direction &
Not a negative kinetic sign; no all-sector energy theorem \\
Tensor and vector sectors in equal-Newton family & Positive listed weights,
real registered speeds & Vacuum/registered background controls, not PPN or
interacting-domain proof \\
Nonzero scalar principal roots & Real with positive squared speeds &
Formal high-momentum symbol only \\
Scalar zero branch & Defective double root & Strong hyperbolicity and
subleading evolution unresolved \\
Fourth-spatial-derivative force branch & Formal \(p^4\) fast mode &
No physical cutoff or scattering window established \\
Constraint singularities & Excluded rather than inverted &
\(k=0\), \(H=0\), zero-density and rank-changing endpoints need separate charts \\
Quantum and renormalized stability & Open & No all-sector determinant or
state-selected semiclassical action \\
\end{tabular}
\end{ruledtabular}
\end{table*}

The correct present classification is therefore
\emph{conditional action with substantive physical holds}.  Check counts are
reproducibility evidence, not a physics pass.

% =============================================================================
\section{Thermodynamic and phenomenological boundaries}
\label{sec:phenomenology}

The reservoir field in Eq.~\eqref{eq:action} is reversible.  For this paper
that is an explicit scope assumption, not an implicit microscopic heat bath:
no temperature, chemical potential, particle-production rate, dissipative
response matrix, or coarse-grained entropy current is specified.  The exchange
currents therefore establish neither positive entropy production nor a
nonzero condensate-number source, and no claim of thermodynamic stability or
an irreversible cosmological arrow follows from this action alone.  A
companion entropy-accounting manuscript may supply the additional constitutive
data required for those separate claims; it is not silently imported here.

Likewise, the current action does not yet provide a validated map to
\(H(z)\), \(G_{\rm eff}(k,z)\), gravitational slip, BBN, a lensing potential,
or a universal galactic force law.  Table~\ref{tab:falsifiers} states
prospective rejection or suspension conditions for this conditional branch;
they are not claims that the listed tests have already been completed.

\begin{table}[t]
\caption{Boundary conditions and future falsifiers for the conditional branch.}
\label{tab:falsifiers}
\begin{ruledtabular}
\begin{tabular}{p{0.46\columnwidth}p{0.46\columnwidth}}
Claim route & Reject or suspend if \\
\hline
Source-vector closure & Direct variation fails total on-shell conservation or
a current becomes singular on the declared finite-field domain \\
Physical scalar EFT & A reduced ghost, nonreal characteristic, unresolved
defective principal block, or nonunique evolution persists on the claimed
domain \\
Weak-field square-root regime & No common regulator-subleading,
quasistatic, periodic and matter-frame domain exists \\
GR recovery & No continuous finite-density path reaches the required local,
cosmological and wave-sector GR observables \\
Cosmological application & One common parameter set fails background,
perturbation, BBN or local-gravity likelihoods \\
\end{tabular}
\end{ruledtabular}
\end{table}

% =============================================================================
\section{Discussion and conclusions}
\label{sec:conclusion}

R4C1 demonstrates that the source-vector problem can be addressed without
postulating \(Q^\nu\) by hand.  A complete declared action yields finite
matter--plenum and plenum--reservoir currents, total on-shell conservation, a
finite-charge background, and the standard uncoupled dust-density kinetic
normalization.

The same calculation narrows the physical claim.  The background cannot
identify the nonlinear spatial-force coefficient; the compact zero mode
forbids a static positive-mean source in the frozen reduction; the universal
pointwise square-root shortcut fails an integrability test; and finite \(b\)
changes the arbitrarily weak-source scaling.  The positive density kinetic
coefficient does not remove the indefinite growing direction or the defective
zero-speed scalar branch.

Accordingly, this paper presents a conditional action, conservation theorem,
normalization control, and scoped obstructions.  It does not claim that ITSM
is established, that \(a_0\) or a \(2/3\) coefficient has been derived, that
all perturbations are stable, or that cosmological data prefer this branch.
Deferred independent review remains required for canonical promotion, while
the scientific holds remain substantive irrespective of review status.

\section*{Code, data, and provenance}

Contracts, derivations, executables, and immutable JSON receipts are available
in the companion repository.  Principal source SHA-256 prefixes are:
action freeze \texttt{81bfc2e4f17b}, first variation
\texttt{4834ea7517fd}, interacting background \texttt{48b49fddd906},
coefficient identifiability \texttt{c4af4084ae83}, density action
\texttt{8a1c2245fc02}, equal-Newton family \texttt{1b9fbc4185b6}, and
equal-Newton scalar analysis \texttt{4d2ea777b8ed}.  Hashes establish artifact
identity, not independent physical validation.

\section*{Declaration on the use of artificial-intelligence tools}

Artificial-intelligence-assisted tools were used for editorial drafting,
source-code assistance, and checks of symbolic and numerical consistency.
The author independently reviewed the manuscript, equations, cited sources,
and generated artifacts and accepts full responsibility for the content.

\begin{acknowledgments}
The author thanks collaborators and reviewers who insisted that action-level
closure, dynamical stability, and phenomenological success be reported as
separate questions.  Remaining errors are the author's.
\end{acknowledgments}

\bibliography{references}

\end{document}
